\section{The compactness alternative and its precise scope}
\label{sec:alternative-audit}

This section reconstructs the finite-dimensional part of
\cite[Theorem D.2, pp.~30--31]{Anthropic2026}.
The inputs are the price-space embedding, the quadrilateral
inequality, and the flat-rectangle statement.
Section~\ref{sec:curvature} explains the curvature argument,
and Lemma~\ref{lem:rectangle} proves the rectangle statement
in the form needed here.
No check of these passages amounts to an independent verification
of either complete circuit-lower-bound preprint.

For matrices $a,W\in\operatorname{Mat}_{\mu}(\F_2)$ and a price
$d$, let
\[
 D(r)^2=\mathbb E_{a,\ \rk W=r}d(a,a+W)^2.
\]
Matrices $W,W'$ are transverse when
$\rk(W+W')=\rk W+\rk W'$.
The same terminology for a list means additivity of the rank
of its sum.
Left and right changes of basis act transitively on transverse
pairs of prescribed ranks, as is seen by putting the two
image spaces and the two row spaces in complementary
coordinate positions.
Their sum is uniform among matrices of the summed rank.
Average the CAT(0) squared-quadrilateral inequality on
$a,a+W,a+W+W',a+W'$.
The two diagonals have the same distribution, giving
\[
 D(i+j)^2\leq D(i)^2+D(j)^2.
\]

\subsection{Compactness has a uniform bound}

Assume the contracting hypothesis fails and fix $\mu\geq3$.
For each $s<\sqrt\mu$ sufficiently close to $\sqrt\mu$,
there is a price with $D(\mu)>sD(1)$.
Here $D(1)>0$: rank-one increments connect the entire
matrix space, so $D(1)=0$ would make every distance zero.
Rescale to $D(1)=1$.

There are $2^{\mu^2}$ matrices and $(2^\mu-1)^2$ nonzero
rank-one matrices. Since the mean square rank-one increment
is one, every such increment has length at most
\[
 B_\mu=\sqrt{2^{\mu^2}(2^\mu-1)^2}.
\]
Every difference is a sum of at most $\mu$ rank-one matrices,
so the whole diameter is bounded by $\mu B_\mu$.
The space of distance arrays is finite-dimensional.
Taking a convergent subsequence as $s\uparrow\sqrt\mu$ is
therefore justified, and its limit is still a price.
Subadditivity gives
\[
 D(r)^2=r \qquad(0\leq r\leq\mu).
\]
Indeed $D(r)^2\leq r$, whereas
$\mu\leq D(\mu)^2\leq D(r)^2+D(\mu-r)^2\leq\mu$.

The averaged quadrilateral inequalities now have zero average
slack. Each slack is nonnegative and only finitely many occur,
so each is zero individually.
Write $n_a(W)=d(a,a+W)^2$.
Equality in the quadrilateral comparison makes opposite sides
equal and gives, for transverse $W,W'$,
\begin{align}
 n_{a+W'}(W)&=n_a(W),\label{eq:audit-side}\\
 d(a+W,a+W')^2
 &=2n_a(W)+2n_a(W')-n_a(W+W').
 \label{eq:audit-diagonal}
\end{align}

\subsection{The exchange argument}

Put
$B_a(W,W')=(n_a(W+W')-n_a(W)-n_a(W'))/2$.
For a transverse triple $W_1,W_2,W_3$,
apply \eqref{eq:audit-diagonal} at $a+W_2$ to
$W_1,W_2+W_3$ and rewrite its terms at $a$.
After cancellation the result is
\[
 n_a(W_1+W_2+W_3)
 -\sum_{i<j}n_a(W_i+W_j)+\sum_i n_a(W_i)=0.
\]
Thus $B_a$ is additive on transverse sums in either argument.

Fix a rank-one $W$ and set
$\chi(u,x)=B_a(W,ux^{\mathsf T})$ when $ux^{\mathsf T}$
is transverse to $W$.
The identity
\[
 u_1x_1^{\mathsf T}+u_2x_2^{\mathsf T}
 =u_1(x_1+x_2)^{\mathsf T}
       +(u_1+u_2)x_2^{\mathsf T}
\]
and the preceding additivity imply
\[
 \chi(u_1,x_1)-\chi(u_1,x_1+x_2)
 =\chi(u_1+u_2,x_2)-\chi(u_2,x_2)
\]
whenever the three matrices $W,u_1x_1^{\mathsf T},u_2x_2^{\mathsf T}$
are transverse. The exchanged pair is then transverse with $W$
as well: its two column vectors span the same space, and so do
its two row vectors.
Replace $x_1$ by $x_1+x_2$.
The left side changes sign and the right side does not;
both vanish over the real numbers.
The coefficient field of the tensors is irrelevant to this
real equality.

For clarity, the connectivity used here is elementary.
Modulo the row space of $W$, the nonzero vectors form a
space of dimension $\mu-1\geq2$.
Vectors in different nonzero quotient classes can be joined
by adding a vector independent of their current class.
Two vectors in the same class can be joined through a
third nonzero class. A complementary column vector can
always be chosen because $\mu\geq3$.
The identical argument changes $u$.
Thus $B_a(W,W')$ is constant as the transverse rank-one
matrix $W'$ varies.
Any two rank-one matrices have a rank-one matrix transverse
to both: choose a column outside their two column lines and
a row outside their two row lines.
Symmetry of $B_a$ consequently gives a common value
$\lambda(a)$ on all transverse rank-one pairs.

Equations \eqref{eq:audit-side}--\eqref{eq:audit-diagonal}
also give
$B_{a+V}(V,W')=-B_a(V,W')$.
For a rank-one $V$ this says
$\lambda(a+V)=-\lambda(a)$.
For rank-two $V$ and a transverse rank-one $W'$, transverse
additivity gives the same sign change, since both $B$ values
are twice the respective $\lambda$.
But rank-two $V$ is also reached by two rank-one increments,
which give no sign change.
Therefore $\lambda(a)=0$ for every $a$.

It follows that $n_a$ is additive on transverse sums.
Apply the same exchange calculation without $W$ to $n_a$
itself. It is constant on rank-one matrices.
Equation \eqref{eq:audit-side} makes that constant independent
of $a$; alternatively one can compare with a rank-one
increment transverse to any prescribed step.
The normalization $D(1)=1$ makes it one.
Decomposing a rank-$r$ matrix into a transverse sum of
rank-one matrices proves
\[
 d(a,a')^2=\rk(a-a').
\]

\subsection{From vertices to a cube}

The diagonal zero-one matrices therefore form an isometric
copy of the vertices of $[0,1]^\mu$.
Fill one coordinate at a time by geodesic segments in the
price space.
At each stage, any two new parallel segments have equal
length one, their two endpoint separations agree, and their
squared diagonal lengths are that common squared separation
plus one.
Lemma~\ref{lem:rectangle} proves that distances between all
points on them are Euclidean.
Induction fills an isometric solid cube.
Every finite subset of $\R^\mu$, after a translation and a
positive rescaling, lies in that cube.
As $\mu\geq3$ was arbitrary, all finite Euclidean metrics are
prices.

The result proved by this argument is precisely
\[
 \text{failure of contraction}\ \Longrightarrow\ \text{flatness}.
\]
It does not prove that contraction excludes flatness.
The source makes the same distinction explicitly in
\cite[Section D.3, p.~31]{Anthropic2026}.
Accordingly the implication from strict shortening to
contraction in Theorem~\ref{thm:shortening-nonflat} is valid,
while its converse is left open.


\subsection{A constant in the partial-transpose comparison}

The factor-two comparison in Corollary~\ref{cor:square-hurwitz}
can also be used without symmetry assumptions on the original
matrix in the direction from partial-transpose rank to sums of
squares. The factor-one inequality for arbitrary matrices printed
in \cite[Lemma~3.12, p.~113:9]{RossmanZhu2026} needs a qualification.
Take $M=E_{(1,1),(2,2)}$, the four-by-four matrix with its only
nonzero entry equal to one in that position. It has rank one
and therefore partial-transpose rank one. Its associated
biquadratic polynomial is $x_1x_2y_1y_2$.
This polynomial is not a square in
$\C[x_1,x_2,y_1,y_2]$, since each of its four irreducible
variable factors has odd multiplicity. Its sum-of-squares
complexity is two, by the product identity in the proof of
Corollary~\ref{cor:square-hurwitz}, with
$\ell=x_1y_1$ and $m=x_2y_2$.
Thus the arbitrary-matrix factor-one statement fails in this
example, whereas the factor-two inequality holds.
The reverse comparison for fully symmetric matrices in
\cite[Lemma~3.14, p.~113:10]{RossmanZhu2026} has the needed
polarization factor; it is proved directly for the crossing
matrix above.
